%% ITEM 0 TYPE Section
\section{Methods}%% ITEM 1 TYPE Section
%% ITEM 1 TYPE Section
\section{Methods: CSPG Approach}%% ITEM 2 TYPE Frame
%% ITEM 2 TYPE Frame
\begin{frame}
\begin{center}
\LARGE\textbf{CSPG: Cross-Partition Sparse Proximity Graph}
\end{center}

\vspace{0.5cm}

\begin{itemize}
    \item \textbf{Goal}: Improve graph-based Approximate Nearest Neighbor Search (ANNS)
    \item \textbf{Core Idea}: Partition dataset to reduce distance computations
    \item \textbf{Key Insight}: Use different step sizes for speed vs. accuracy
    \item \textbf{Two Key Algorithms}: Partitioning algorithm and search algorithm
    \item \textbf{Approach}: Balance between fast approaching and precise search
\end{itemize}
\end{frame}%% ITEM 3 TYPE Frame
%% ITEM 3 TYPE Frame
\begin{frame}
\begin{center}
\Large\textbf{Partitioning Algorithm: Core Strategy}
\end{center}

\vspace{0.3cm}

\begin{itemize}
    \item \textbf{Random Partitioning}: Dataset $\mathcal{D}$ divided into multiple groups
    \item \textbf{Routing Vectors}: Common vectors shared across groups (red nodes)
    \item \textbf{Group-Specific Vectors}: Distinct vectors in each group (green/blue nodes)
    \item \textbf{Sparse Subgraphs}: Build smaller, sparser graphs for each group
    \item \textbf{Key Benefit}: Sparse graphs allow larger steps for fast approaching
\end{itemize}

\vspace{0.3cm}

\textbf{Example}: Two groups $\mathcal{G}_1$ (left) and $\mathcal{G}_2$ (right) with shared routing vectors
\end{frame}%% ITEM 4 TYPE Frame
%% ITEM 4 TYPE Frame
\begin{frame}
\begin{center}
\Large\textbf{Search Algorithm: Two-Stage Approach}
\end{center}

\vspace{0.3cm}

\begin{itemize}
    \item \textbf{Stage 1}: Fast approaching using single partition
    \item \textbf{Stage 2}: Precise search across all partitions
    \item \textbf{Different Beam Sizes}: $ef_1 < ef_2$ for two stages
    \item \textbf{Traditional vs CSPG}: Fixed candidate set vs dynamic expansion
    \item \textbf{Goal}: Quickly zone in, then carefully search around
\end{itemize}

\vspace{0.3cm}

\textbf{Key Innovation}: Dynamic traversal across multiple proximity graphs
\end{frame}%% ITEM 5 TYPE Frame
%% ITEM 5 TYPE Frame
\begin{frame}
\begin{center}
\Large\textbf{Stage 1: Fast Approaching}
\end{center}

\vspace{0.3cm}

\begin{itemize}
    \item \textbf{Greedy Beam Search}: Starts from group $\mathcal{G}_1$
    \item \textbf{Short Search Length}: Fewer neighbor expansions
    \item \textbf{Visualization}: Yellow arrows in method diagram
    \item \textbf{Result}: Zones in on query's nearby region
    \item \textbf{Output}: Candidate set with closest vector from first partition
\end{itemize}

\vspace{0.3cm}

\textbf{Process}: Search from $\mathcal{G}_1$ → Approach query → Continue from $\mathcal{G}_2$
\end{frame}%% ITEM 6 TYPE Frame
%% ITEM 6 TYPE Frame
\begin{frame}
\begin{center}
\Large\textbf{Stage 2: Cross-Partition Expansion}
\end{center}

\vspace{0.3cm}

\begin{itemize}
    \item \textbf{Dynamic Graph Traversal}: Search across multiple proximity graphs
    \item \textbf{Routing Vector Expansion}: If neighbor $u$ is routing vector, add all instances
    \item \textbf{Maximized Search Space}: Include all potential closest vectors
    \item \textbf{Visualization}: Orange arrows in method diagram
    \item \textbf{Key Feature}: Appears as single graph search in practice
\end{itemize}

\vspace{0.3cm}

\textbf{Critical Difference}: Routing vectors trigger cross-partition expansion
\end{frame}%% ITEM 7 TYPE Frame
%% ITEM 7 TYPE Frame
\begin{frame}
\begin{center}
\Large\textbf{CSPG Advantages and Performance}
\end{center}

\vspace{0.3cm}

\begin{itemize}
    \item \textbf{Speed Advantage}: Sparse graphs enable faster initial search
    \item \textbf{Accuracy Advantage}: Cross-partition expansion ensures thorough search
    \item \textbf{Balance Achieved}: Large steps when far, small steps when close
    \item \textbf{Outperforms Traditional PG}: With appropriate partitioning
    \item \textbf{Generalizable}: Works with multiple partitions beyond two
\end{itemize}

\vspace{0.3cm}

\textbf{Key Result}: Higher query speed and accuracy compared to traditional proximity graphs
\end{frame}%% ITEM 8 TYPE Frame
%% ITEM 8 TYPE Frame
\begin{frame}
\begin{center}
\Large\textbf{Summary: CSPG Method}
\end{center}

\vspace{0.3cm}

\begin{itemize}
    \item \textbf{Partitioning}: Random division with shared routing vectors
    \item \textbf{Two-Stage Search}: Fast approaching + cross-partition expansion
    \item \textbf{Speed Mechanism}: Sparse graphs enable large steps when far
    \item \textbf{Accuracy Mechanism}: Routing vectors enable thorough search
    \item \textbf{Key Innovation}: Dynamic traversal across multiple graphs
\end{itemize}

\vspace{0.3cm}

\textbf{Final Result}: Balanced approach outperforming traditional proximity graphs
\end{frame}%% ITEM 9 TYPE Section
%% ITEM 9 TYPE Section
\section{Conclusion}%% ITEM 10 TYPE Section
%% ITEM 10 TYPE Section
\section{Conclusion: CSPG Performance Summary}%% ITEM 11 TYPE Frame
%% ITEM 11 TYPE Frame
\begin{frame}
\begin{center}
\LARGE\textbf{Conclusion: CSPG Performance Summary}
\end{center}

\vspace{10pt}

\textbf{CSPG achieves superior performance through innovative design:}

\begin{itemize}
\item \textbf{Higher QPS and Recall@10} than traditional proximity graphs
\item \textbf{Random partitioning} creates smaller, sparser graphs per partition
\item \textbf{Efficient two-stage search} with routing vectors for cross-partition expansion
\item \textbf{Consistent improvements} across multiple baseline algorithms
\item \textbf{Balanced approach} for both speed and accuracy
\end{itemize}
\end{frame}%% ITEM 12 TYPE Frame
%% ITEM 12 TYPE Frame
\begin{frame}
\textbf{Core Innovation: Random Partitioning}

\begin{itemize}
\item \textbf{Dataset $\mathcal{D}$ is randomly divided} into multiple groups
\item \textbf{Routing vectors} are shared across groups for connectivity
\item \textbf{Smaller, sparser graphs} are built for each partition
\item \textbf{Enables larger search steps} when far from query (for speed)
\item \textbf{Allows finer search} when close to query (for accuracy)
\end{itemize}

\vspace{10pt}

\textbf{Key Benefits:}
\begin{itemize}
\item Reduced graph complexity per partition
\item Faster initial approach to query region
\item Maintains accuracy through routing connections
\end{itemize}
\end{frame}%% ITEM 13 TYPE Frame
%% ITEM 13 TYPE Frame
\begin{frame}
\textbf{Efficient Two-Stage Search Algorithm}

\begin{itemize}
\item \textbf{First Stage: Fast Approaching}
\begin{itemize}
\item Greedy beam search within a single partition
\item Quickly approaches the query region
\item Uses sparser graph for larger steps
\end{itemize}

\item \textbf{Second Stage: Precise Search}
\begin{itemize}
\item Cross-partition expansion using routing vectors
\item When encountering a routing vector, all its instances across partitions are added to candidate set
\item Maximizes search space for accurate results
\end{itemize}

\item \textbf{Smart Transition}
\begin{itemize}
\item Automatically switches between stages
\item Balances speed and accuracy dynamically
\end{itemize}
\end{itemize}
\end{frame}%% ITEM 14 TYPE Frame
%% ITEM 14 TYPE Frame
\begin{frame}
\textbf{Experimental Results: Superior Performance}

\begin{itemize}
\item \textbf{CSPG built upon multiple baselines:}
\begin{itemize}
\item HNSW (Hierarchical Navigable Small World)
\item HCNNG (Hierarchical Clustering Nearest Neighbor Graph)
\item Vamana
\item NSG (Navigating Spreading-out Graph)
\end{itemize}

\item \textbf{Consistent improvements across all baselines}
\item \textbf{Higher QPS:} More queries processed per second
\item \textbf{Higher Recall@10:} More accurate nearest neighbor retrieval
\item \textbf{Better trade-off curve} between speed and accuracy
\end{itemize}

\vspace{10pt}

\textbf{Key Finding:} CSPG's design principles work effectively regardless of the underlying graph construction algorithm.
\end{frame}%% ITEM 15 TYPE Frame
%% ITEM 15 TYPE Frame
\begin{frame}
\textbf{Why CSPG Works: Key Insights}

\begin{itemize}
\item \textbf{Smaller graphs are faster to search}
\begin{itemize}
\item Reduced complexity per partition
\item Faster neighbor exploration
\end{itemize}

\item \textbf{Routing vectors enable efficient cross-partition search}
\begin{itemize}
\item Maintain connectivity between partitions
\item Enable smart expansion when needed
\end{itemize}

\item \textbf{Two-stage approach matches search strategy to proximity}
\begin{itemize}
\item Fast when far, precise when close
\item Dynamic adaptation to query location
\end{itemize}

\item \textbf{Random partitioning ensures balanced distribution}
\begin{itemize}
\item Avoids bias in partition creation
\item Maintains statistical properties
\end{itemize}
\end{itemize}
\end{frame}%% ITEM 16 TYPE Frame
%% ITEM 16 TYPE Frame
\begin{frame}
\begin{center}
\LARGE\textbf{Final Summary}
\end{center}

\vspace{15pt}

\textbf{CSPG delivers:}

\begin{itemize}
\item \textbf{Superior performance} in both speed (QPS) and accuracy (Recall@10)
\item \textbf{Innovative architecture} combining random partitioning with two-stage search
\item \textbf{Consistent improvements} across multiple graph-based ANNS algorithms
\item \textbf{Efficient cross-partition exploration} through routing vectors
\item \textbf{Practical solution} for large-scale nearest neighbor search problems
\end{itemize}

\vspace{15pt}

\textbf{Key Contribution:} A novel approach that fundamentally improves the speed-accuracy trade-off in proximity graph search.
\end{frame}%% ITEM 17 TYPE Frame
%% ITEM 17 TYPE Frame
\begin{frame}
\begin{center}
\LARGE\textbf{Thank You}
\end{center}

\vspace{20pt}

\begin{center}
\textbf{Questions?}
\end{center}

\vspace{20pt}

\begin{center}
\textbf{Key Takeaways:}
\end{center}

\begin{itemize}
\item CSPG achieves higher QPS and Recall@10
\item Random partitioning creates efficient smaller graphs
\item Two-stage search with routing vectors enables cross-partition expansion
\item Consistent improvements across multiple baseline algorithms
\item Better speed-accuracy trade-off for ANNS
\end{itemize}
\end{frame}